% Custom CV — Fraunhofer IGB Stuttgart — Master Thesis: green reactor for wastewater treatment (52683) — 2026-11-12
% Base: resumes/library/electrochem_en/source.tex (no EN water CV in the library; water projects from me.md P9/P10)
\documentclass[10pt]{article}
\usepackage[a4]{nermine-cv}
\begin{document}
\cvheader{Nermine Moumen}{Final-year Engineering Student in Industrial Chemistry}{nermine.moumen@insat.ucar.tn}{+21627040414}{linkedin.com/in/nermine-moumen}

\section{EDUCATION}
\entry{National Institute of Applied Sciences and Technology (INSAT)}{\textit{2022–Present}}
\summary{Industrial Chemical Engineering student, specialized in Process and Environment, with a strong foundation in water treatment, analytical chemistry and electrochemical processes. Oriented towards laboratory-scale experiments, process monitoring and data interpretation. Seeking a 6-month master thesis internship starting February 2027.}

\section{PROFESSIONAL EXPERIENCE}
\entry{SEGULA Technologies – Electrochemical Battery Modeling Internship}{\textit{06/2026–08/2026}}
\begin{bul}
\item Modeled a \textbf{Li-ion cell} (NMC811/Graphite-SiOx) using the Single Particle Model (SPM) implemented in MATLAB.
\item Implemented the electrodes' \textbf{open-circuit potential (OCP) curves} and solved the PDEs with \textbf{FVM}.
\item Optimized the mesh and diffusion coefficient values, reducing the RMSE error \textbf{from 202.13 mV to 40.5 mV}.
\item Extracted and processed \textbf{experimental discharge data} from a database to validate MATLAB and COMSOL results.
\end{bul}
\entry{Control Union – Quality Control Internship}{\textit{08/2025–09/2025}}
\begin{bul}
\item Performed \textbf{laboratory analysis} of lubricants (viscosity, heavy metals by \textbf{ICP-OES}) following standardized procedures.
\item Prepared and interpreted \textbf{analytical reports} to verify compliance with ASTM D5185 and \textbf{ISO 17025} standards.
\end{bul}
\entry{ETAP – Water Treatment and Petrochemistry Internship}{\textit{07/2024–08/2024}}
\begin{bul}
\item Performed physico-chemical and ionic analyses of water samples (pH, salinity, chlorides via ion chromatography).
\item Studied water treatment methods \textbf{(coagulation-flocculation, ion exchange)} for contaminant and \textbf{corrosion} control.
\item Performed organic geochemistry techniques \textbf{(GC, GC-MS, Rock-Eval pyrolysis)} for hydrocarbon \textbf{characterization}.
\end{bul}
\section{ADDITIONAL EXPERIENCE}
\entry{Cascade – Online Chemistry Teacher (Part-time)}{\textit{2024–Present}}
\begin{bul}
\item Provide individualized online chemistry tutoring for French Terminale students, covering exam-level concepts.
\end{bul}

\section{ACADEMIC PROJECTS}
\project{Electrochemical Treatment of an Industrial Effluent by Electroflocculation}
\begin{bul}
\item Carried out \textbf{electrochemical measurements} on aluminum \textbf{electrodes}: current-potential curves and \textbf{pulse polarography}.
\item Calculated the \textbf{faradaic yield} with \textbf{Faraday's law}, linking electrode dissolution to the charge passed.
\item Removal of Cd$^{2+}$ and Zn$^{2+}$ from a simulated effluent: \textbf{85\% and 57\%} in 20 min.
\end{bul}
\project{Effluent Treatment by Coagulation-Flocculation}
\begin{bul}
\item Optimized the coagulation pH (aluminium sulfate / alginic acid) on a \textbf{real effluent}, testing 5 pH values; maximum settled floc volume at \textbf{pH 8}.
\item Interpreted the efficiency through aluminium speciation using the \textbf{Pourbaix diagram}.
\end{bul}
\project{Energy Recovery from Organic Waste by Anaerobic Digestion (ASBR)}
\begin{bul}
\item Experimental follow-up of \textbf{5 treatment configurations} (primary, secondary, mixed sludge and co-digestion with fruit and vegetable waste) in \textbf{ASBR digesters} (HRT = 20 d).
\item Compared performance through COD, VSS, volumetric load, biogas production and CH$_4$ content, reaching up to \textbf{0.50 L/d of biogas at 64.52\% CH$_4$}.
\end{bul}

\section{TECHNICAL SKILLS}
\begin{bul}
\item \textbf{Lab Techniques:} Chromatography (HPLC, GC, TLC), Spectrometry (MS, ICP-OES), Spectroscopy (NMR, FTIR, UV-Vis), Diffraction (XRD), Microscopy (SEM, TEM), Electrochemistry (Cyclic voltammetry, Polarography).
\item \textbf{Simulation:} Aspen Plus, Aspen HYSYS, COMSOL Multiphysics\sep\textbf{Programming:} \textit{MATLAB, Mathematica, Python}
\item \textbf{Data \& Office:} Microsoft Office Suite, OpenLCA\sep\textbf{Standards:} ISO 9001, ISO 14001, ISO 45001, ISO 17025, ISO 14040
\end{bul}

\section{EXTRACURRICULAR ACTIVITIES}
\begin{bul}
\item \textbf{IEEE INSAT} – Chair of EMBS; Organizing Chair of TRSYP 3.0; Social Activities Chair of MENASYP.
\item \textbf{THEATRO INSAT} – Member of the dance team.
\end{bul}

\section{LANGUAGES AND CERTIFICATES}
\twocol{Languages}{English (Fluent), French (Fluent), Arabic (Native).}{Certificates}{Aspen HYSYS – MATLAB}
\end{document}
